\begin{table}[htbp]
\centering
\small
\begin{tabularx}{\textwidth}{|c|>{\centering\arraybackslash}X|c|c|c|c|>{\centering\arraybackslash}p{2.8cm}|}
\hline
\textbf{№ эксп.} & \textbf{\begin{tabular}[c]{@{}c@{}}№ работающего ПЧ /\\ Срабатывание БАВР\end{tabular}} & \textbf{\begin{tabular}[c]{@{}c@{}}Ввод 1\\ (Т1)\end{tabular}} & \textbf{\begin{tabular}[c]{@{}c@{}}Ввод 2\\ (Т2)\end{tabular}} & \textbf{$i_{\max}$} & \textbf{$K_{\text{броск}}$} & \textbf{\commentnote{DISS-005-C0055}{Feel}{Вот это ч хороший столбец\newline А вывод-то какой?}{\commentnote{DISS-005-C0056}{Евдаков Алексей}{Никакой)) Я прост вставил таблицу, в которой у меня были указаны эти колебания в разных опытах. Там анализировался факт отключения нагрузки при БАВР, и то, что дело не в скорости работы, колебания и т.д. А в том, в какой момент этого колебания прост случилась авария (отключение выключателя).}{Время (от начала аварии до команд БАВР)}}} \\ \hline
\multicolumn{7}{|c|}{\textbf{Отключения ввода №2 (Т2)}} \\ \hline
\multirow{3}{*}{1} & № работающего ПЧ & 1 & 5 & \multirow{3}{*}{491~А} & \multirow{3}{*}{4,6} & \multirow{3}{*}{31,667~мс} \\ \cline{2-4}
 & Срабатывание БАВР & + & -- & & & \\ \cline{2-4}
 & Уровень колебаний & 5\% & 34\% & & & \\ \hline
\multirow{3}{*}{2} & № работающего ПЧ & 1 & 6 & \multirow{3}{*}{305~А} & \multirow{3}{*}{2,3} & \multirow{3}{*}{33,333~мс} \\ \cline{2-4}
 & Срабатывание БАВР & + & + & & & \\ \cline{2-4}
 & Уровень колебаний & 10\% & 38\% & & & \\ \hline
\multirow{3}{*}{3} & № работающего ПЧ & 3 & 5 & \multirow{3}{*}{320~А} & \multirow{3}{*}{3} & \multirow{3}{*}{34,167~мс} \\ \cline{2-4}
 & С Срабатывание БАВР & + & + & & & \\ \cline{2-4}
 & Уровень колебаний & 8\% & 34\% & & & \\ \hline
\multicolumn{7}{|c|}{\textbf{Отключения ввода №1 (Т1)}} \\ \hline
\multirow{3}{*}{4} & № работающего ПЧ & 1 & 6 & \multirow{3}{*}{270~А} & \multirow{3}{*}{2,3} & \multirow{3}{*}{17,5~мс} \\ \cline{2-4}
 & Срабатывание БАВР & -- & + & & & \\ \cline{2-4}
 & Уровень колебаний & 28\% & 31\% & & & \\ \hline
\multirow{3}{*}{5} & № работающего ПЧ & 2 & 5 & \multirow{3}{*}{350~А} & \multirow{3}{*}{3,2} & \multirow{3}{*}{29,167~мс} \\ \cline{2-4}
 & Срабатывание БАВР & + & + & & & \\ \cline{2-4}
 & Уровень колебаний & 36\% & 22\% & & & \\ \hline
\multirow{3}{*}{6} & № работающего ПЧ & 3 & 5 & \multirow{3}{*}{242~А} & \multirow{3}{*}{2,2} & \multirow{3}{*}{30,833~мс} \\ \cline{2-4}
 & Срабатывание БАВР & + & + & & & \\ \cline{2-4}
 & Уровень колебаний & 17\% & 40\% & & & \\ \hline
\end{tabularx}
\end{table}